\subsection{Арифметика и сравнения}

Операции: \verb|+ - * / %| (остаток от деления) и \verb|^| (степень). Сравнения: \verb|< <= > >= ==| и \verb|~=| («не равно»). Приоритет как в математике: сначала \verb|^|, затем \verb|* / %|, потом \verb|+ -|; скобки меняют порядок.

\begin{codefile}{lua/04_language/lua_arithmetic_ops.lua}
\end{codefile}

\paragraph{Что происходит.}
\begin{itemize}
    \item \verb|2 ^ 3| даёт \verb|8.0|, а не \verb|8|: степень и деление \verb|/| всегда возвращают вещественное число. Целые и вещественные числа в Lua~5.3 --- два подвида одного типа, в арифметике они смешиваются свободно.
    \item Оператора \verb|**| в Lua нет --- это Python. Степень только через \verb|^|.
    \item \verb|math.floor(17 / 4)| --- целочисленное деление с округлением вниз. В симуляторе работает и \verb|//|, но \verb|math.floor| переносимее.
    \item \verb|math.sqrt|, \verb|math.abs|, \verb|math.max|, \verb|math.min|, \verb|math.pi| --- то, что понадобится в миссиях. Тригонометрия Lua и автопилот измеряют углы в \textbf{радианах}: полный оборот --- $2\pi \approx 6.28$~рад, $180^\circ = \pi$. Перевод градусов в радианы: \verb|math.rad(deg)|, обратно --- \verb|math.deg(rad)|.
\end{itemize}

\begin{hintbox}[Ограничить значение диапазоном]
\verb|v = math.max(0, math.min(1, v))| «зажимает» \verb|v| между 0 и 1. Этот приём пригодится для яркости светодиодов.
\end{hintbox}

\subsubsection*{Задачи для самостоятельного решения}
\begin{enumerate}
  \item Напечатайте \verb|3 + 4 * 2 - 5| и \verb|(3 + 4) * (2 - 5)|.
  \item Напечатайте остаток \verb|17 % 4| и проверьте условием, чётно ли число \verb|n = 10|.
  \item Вычислите длину вектора \verb|(x, y) = (3, 4)| через \verb|math.sqrt(x*x + y*y)|.
  \item Переведите \verb|90| градусов в радианы двумя способами: \verb|deg * math.pi / 180| и \verb|math.rad(deg)|. Сравните результаты.
  \item Реализуйте линейную интерполяцию \verb|lerp = a + t * (b - a)| для \verb|a = 0|, \verb|b = 10|, \verb|t = 0.25|.
  \item Напечатайте среднее трёх чисел \verb|(a + b + c) / 3|.
\end{enumerate}
